% ============================================================================
% Capítulo 6. Resultados
% Borrador: sección 6.1 redactada; secciones 6.2 a 6.5 pendientes.
%
% Las tablas se insertan desde analisis/tablas/latex y las figuras desde
% analisis/figuras. Si en el proyecto de LaTeX esas carpetas tienen otra ruta,
% basta cambiar las dos líneas siguientes (o definir los comandos antes de
% incluir este archivo).
% Requiere: booktabs, graphicx, amsmath, amssymb.
%
% Las cifras del texto que no están en una tabla se pueden volver a calcular
% con texto/cifras_seccion_6_1.py.
% ============================================================================
\providecommand{\rutatablas}{tablas/latex}
\providecommand{\rutafiguras}{figuras}

\chapter{Resultados}
\label{cap:resultados}

Este capítulo presenta los resultados en el orden en que se plantearon los análisis.
Primero se describen el catálogo depurado y las submuestras que se obtienen de él
(Sección~\ref{sec:res-muestras}). Después se presentan los tres análisis: los regímenes de
radio (Sección~\ref{sec:res-analisis-I}), las propiedades físicas de los planetas pequeños
(Sección~\ref{sec:res-analisis-II}) y la arquitectura de los sistemas
(Sección~\ref{sec:res-analisis-III}). El capítulo termina con una comparación de la
contribución estelar en los tres (Sección~\ref{sec:res-comparacion}). Las tablas y figuras
complementarias se reúnen en el Anexo~\ref{anx:complementarias}.

% ----------------------------------------------------------------------------
\section{Descripción de las muestras}
\label{sec:res-muestras}

\subsection{Catálogo general y disponibilidad de parámetros}
\label{sec:res-catalogo}

\paragraph{Depuración.}
La tabla descargada contiene 6\,324 planetas confirmados alrededor de 4\,738 estrellas
anfitrionas. No se encontraron planetas duplicados ni valores menores o iguales a cero en
las variables físicas. Se excluyeron los 57 planetas que el archivo marca como
controversiales, de modo que el catálogo depurado tiene 6\,267 planetas y 4\,711 estrellas
(Tabla~\ref{tab:D-T1}). Todos los resultados de este capítulo se refieren a ese catálogo.
De las 4\,711 estrellas, 1\,039 (22.1\,\%) tienen dos o más planetas en el catálogo, y esos
sistemas reúnen 2\,595 planetas (41.4\,\%). Una parte importante de los planetas comparte,
por tanto, estrella anfitriona con otros.

\paragraph{Método de descubrimiento.}
Tres de cada cuatro planetas se descubrieron por tránsito (4\,660; 74.4\,\%). Les siguen
la velocidad radial (1\,162; 18.5\,\%), la microlente gravitacional (281; 4.5\,\%) y la
imagen directa (95; 1.5\,\%); los siete métodos restantes suman 69 planetas (1.1\,\%)
(Tabla~\ref{tab:D-T4} y Figura~\ref{fig:D-F2}). Tres misiones espaciales concentran dos
tercios del catálogo: \textit{Kepler} (2\,774 planetas), TESS (908) y K2 (549), que en
conjunto aportan el 67.5\,\%.

\paragraph{Disponibilidad de parámetros.}
Las variables que usan los modelos están disponibles para casi todo el catálogo: la masa
estelar en el 99.9\,\% de los planetas, el radio planetario en el 99.3\,\%, el período
orbital en el 94.6\,\% y la metalicidad en el 90.0\,\% (Tabla~\ref{tab:D-T2}). La
metalicidad es el predictor principal con más datos faltantes (628 planetas) y es la
variable que más reduce las submuestras. La temperatura efectiva y el radio estelar, que
se necesitan para calcular el flujo incidente, faltan en cerca del 5\,\% de los planetas.
Las variables más incompletas son el tipo espectral (63.3\,\% de datos faltantes), la edad
estelar (22.6\,\%) y la excentricidad (16.6\,\%) (Figura~\ref{fig:D-F1}).

Que un parámetro esté disponible no significa que se haya medido. La tabla de parámetros
compuestos completa el radio o la masa que faltan con una relación masa--radio. De los
6\,223 planetas con radio, 4\,701 (75.5\,\%) se observaron en tránsito; el radio de los
1\,522 restantes (24.5\,\%) no proviene de un tránsito. En cuanto a la masa, el catálogo la
identifica como masa verdadera en 2\,390 planetas (38.1\,\%), como masa mínima,
$M_p \sin i$, en 903 (14.4\,\%) y como estimación a partir del radio en 2\,974 (47.5\,\%);
estos últimos son todos planetas observados en tránsito. Solo una fracción del catálogo
tiene, entonces, radio y masa medidos de forma independiente, y por eso cada análisis
exige la medición que necesita en lugar de usar todos los valores disponibles.

La metalicidad tampoco es homogénea. De los 5\,639 planetas con metalicidad, en 5\,044
(89.4\,\%) el valor corresponde a [Fe/H] y en 578 (10.3\,\%) a [M/H]; en 17 el catálogo no
indica la escala. En este trabajo ambas se trataron como una sola variable.

\paragraph{Estadística descriptiva.}
La Tabla~\ref{tab:D-T3} resume las variables planetarias y estelares principales. Se
reportan la mediana y los cuartiles porque las distribuciones son asimétricas y abarcan
varios órdenes de magnitud. El planeta mediano del catálogo tiene un radio de
2.85~$R_\oplus$ y un período de 10.6~días. La mitad de los períodos está entre 4.3 y
37.6~días; el 48.2\,\% es menor que 10~días y el 83.6\,\%, menor que 100~días. La estrella
anfitriona mediana tiene 0.94~$M_\odot$, 5\,544~K y [Fe/H]~$= 0.02$. El intervalo
intercuartílico de la masa estelar va de 0.77 a 1.08~$M_\odot$ y el de la metalicidad, de
$-0.08$ a 0.13~dex. El 86.9\,\% de los planetas orbita estrellas de entre 0.5 y
1.5~$M_\odot$, y en el 89.4\,\% la metalicidad difiere de la solar en menos de 0.3~dex. El
catálogo describe, por lo tanto, sobre todo planetas de período corto alrededor de
estrellas parecidas al Sol, y el intervalo de masa estelar y de metalicidad que cubre es
estrecho.

\input{\rutatablas/D-T3_descriptiva_variables_principales}

Los valores extremos de la Tabla~\ref{tab:D-T3} no representan a la población. No se aplicó
un corte en masa, y 151 objetos (2.4\,\%) tienen una masa catalogada mayor que 13 masas de
Júpiter; la masa máxima, 9\,535~$M_\oplus$ (30 masas de Júpiter), es la de un compañero
descubierto por imagen directa. La densidad máxima corresponde a un planeta cuya masa es un
límite superior, y el período máximo, a un planeta de imagen directa en una órbita muy
amplia. Estos valores no afectan a los modelos: el análisis~II excluye los límites y se
restringe a planetas de 1 a 4~$R_\oplus$, y los demás análisis no usan la masa planetaria.
La muestra del análisis~I sí incluye 13 objetos (0.3\,\%) con masa medida mayor que 13
masas de Júpiter.

\paragraph{Regímenes de radio.}
De los 6\,223 planetas con radio, el régimen más poblado es el de tipo Neptuno (2\,374;
38.1\,\%), seguido por los gigantes gaseosos (1\,853; 29.8\,\%), las super-Tierras (1\,188;
19.1\,\%), los planetas terrestres (571; 9.2\,\%) y los mayores que Júpiter (237; 3.8\,\%)
(Tabla~\ref{tab:D-T5}). La distribución de radios tiene dos máximos, uno entre 2 y
3~$R_\oplus$ y otro entre 11 y 14~$R_\oplus$, con pocos planetas entre 4 y 10~$R_\oplus$
(Figura~\ref{fig:D-F3}a).

\input{\rutatablas/D-T5_regimenes_de_radio}

\begin{figure}[htbp]
    \centering
    \includegraphics[width=\textwidth]{\rutafiguras/D-F3_distribucion_radio_planetario}
    \caption{Distribución del radio planetario. (a) Todo el catálogo; las líneas verticales
    marcan los límites de los cinco regímenes (1.25, 2, 6 y 15~$R_\oplus$). (b) Planetas en
    tránsito entre 0.5 y 4~$R_\oplus$; la banda gris marca el intervalo
    1.6--2.0~$R_\oplus$ excluido en el análisis de super-Tierras y sub-Neptunos.}
    \label{fig:D-F3}
\end{figure}

El segundo máximo depende de cómo se obtuvo el radio. En los demás regímenes entre el 87 y
el 95\,\% de los planetas se observó en tránsito, pero entre los gigantes gaseosos solo 752
de 1\,853 (40.6\,\%). Los 1\,101 gigantes sin tránsito tienen radios muy concentrados
(cuartiles de 12.4 y 13.6~$R_\oplus$): son los que forman la barra gris más alta de la
Figura~\ref{fig:D-F3}a, y su radio refleja la relación masa--radio del archivo y no una
medición. Por esta razón los análisis~I y~II usan solo planetas de tránsito, y en el
análisis~III los resultados se comprueban en los sistemas cuyos planetas se observaron
todos en tránsito.

Las propiedades estelares cambian con el régimen (Tabla~\ref{tab:D-T5}). La mediana de la
masa estelar aumenta de forma monótona con el tamaño del planeta: 0.83~$M_\odot$ para los
terrestres, 0.89 para las super-Tierras, 0.92 para los de tipo Neptuno, 1.03 para los
gigantes gaseosos y 1.30 para los mayores que Júpiter. La metalicidad mediana es
prácticamente solar en los tres regímenes de menor tamaño ($-0.02$, 0.00 y 0.01~dex) y más
alta en los gigantes gaseosos (0.10~dex) y en los mayores que Júpiter (0.05~dex). El
período mediano también difiere: 4.5, 6.5 y 14.2~días en los tres primeros regímenes, y
85.3~días en los gigantes gaseosos. Este último valor mezcla dos poblaciones: los gigantes
sin tránsito, descubiertos en su mayoría por velocidad radial, tienen un período mediano de
694~días, y los que transitan, de 4.6~días. Los planetas mayores que Júpiter son casi todos
de período corto (213 de 237 tienen $P < 10$~días; mediana de 3.3~días).

Estas diferencias son descriptivas. No tienen en cuenta el período orbital ni el método de
descubrimiento, de modo que no permiten separar una asociación con la estrella de un efecto
de la detección. Esa separación es el objeto del análisis~I
(Sección~\ref{sec:res-analisis-I}).

\paragraph{Planetas pequeños.}
Entre los planetas en tránsito de 0.5 a 4~$R_\oplus$, la distribución de radios tiene dos
máximos separados por un mínimo (Figura~\ref{fig:D-F3}b). En intervalos de
0.1~$R_\oplus$, el primer máximo está en 1.5--1.6~$R_\oplus$, con 178 planetas, y el
segundo en 2.2--2.4~$R_\oplus$, con 167 planetas en cada intervalo; el mínimo, en
1.8--1.9~$R_\oplus$, tiene 111. Este mínimo queda dentro del intervalo de 1.6 a
2.0~$R_\oplus$ que se excluye al comparar super-Tierras y sub-Neptunos. La depresión es
moderada: el intervalo del mínimo contiene el 62\,\% de los planetas que el del primer
máximo.

\paragraph{Diagrama masa--radio.}
La Figura~\ref{fig:D-F4} muestra los 1\,504 planetas (24.0\,\% del catálogo, en 1\,224
estrellas) que se observaron en tránsito y tienen masa verdadera, sin límites en la masa ni
en el radio. Se distinguen tres comportamientos. Por debajo de 1.6~$R_\oplus$ los planetas
se agrupan alrededor de la curva de densidad terrestre: su densidad mediana es de
6.7~g/cm$^{3}$ (cuartiles de 5.1 y 10.3; $N = 141$). Entre 2 y 4~$R_\oplus$ la densidad
mediana baja a 2.8~g/cm$^{3}$ (cuartiles de 1.8 y 4.2; $N = 369$) y, a un mismo radio, las
masas se dispersan casi un orden de magnitud: en el régimen de tipo Neptuno, los
percentiles 10 y 90 de la masa son 4.6 y 39.3~$M_\oplus$. En los gigantes gaseosos el radio
deja de informar sobre la masa: los percentiles 10 y 90 son 50 y 1\,215~$M_\oplus$ dentro
de un intervalo de radios estrecho. La dispersión de masas a radio fijo entre los planetas
pequeños es la variación que examina el análisis~II (Sección~\ref{sec:res-analisis-II}).

\begin{figure}[htbp]
    \centering
    \includegraphics[width=0.82\textwidth]{\rutafiguras/D-F4_diagrama_masa_radio}
    \caption{Diagrama masa--radio de los planetas en tránsito con masa verdadera, por
    régimen de radio. Las curvas grises son líneas de densidad constante (1 y
    5.51~g/cm$^{3}$).}
    \label{fig:D-F4}
\end{figure}

\paragraph{Parámetros derivados.}
El flujo incidente se calculó para 5\,948 planetas, el 94.9\,\% del catálogo y 253 más
de los que tienen un flujo catalogado. En 412 de esos 5\,948 planetas el semieje mayor se
obtuvo con la tercera ley de Kepler. En los planetas que tienen ambos valores, la razón entre el flujo del
catálogo y el calculado tiene una mediana de 0.997 y cuartiles de 0.919 y 1.080
($N = 5\,644$), y la razón entre el semieje mayor del catálogo y el de Kepler, una mediana
de 1.000 y cuartiles de 0.994 y 1.006 ($N = 5\,513$) (Tabla~\ref{tab:D-T7}). El flujo
calculado reproduce, entonces, el del catálogo. Las clasificaciones auxiliares por masa y
por período, que no se usan en los modelos, se resumen en la Tabla~\ref{tab:D-T6}.

% ----------------------------------------------------------------------------
\subsection{Submuestras analíticas}
\label{sec:res-submuestras}

La Tabla~\ref{tab:D-T8} resume las cuatro submuestras. Ninguna se obtuvo con un filtro
único sobre el catálogo: cada una conserva los planetas que tienen completas las variables
de su modelo.

\input{\rutatablas/D-T8_submuestras_analiticas}

\paragraph{Análisis I.}
De los 4\,660 planetas descubiertos por tránsito, 4\,461 (95.7\,\%) tienen radio, período,
masa estelar y metalicidad. De los 199 excluidos, 198 carecen de metalicidad. La muestra
abarca 3\,390 estrellas, de las cuales 706 (20.8\,\%) aportan más de un planeta. El
61.5\,\% de los planetas proviene de \textit{Kepler}, el 17.1\,\% de TESS y el 11.3\,\% de
K2.

Para la comparación entre super-Tierras y sub-Neptunos, el catálogo tiene 3\,261 planetas
descubiertos por tránsito con radio entre 1 y 4~$R_\oplus$: 873 super-Tierras, 568 en el
intervalo de 1.6 a 2.0~$R_\oplus$ y 1\,820 sub-Neptunos. Al excluir ese intervalo y exigir
flujo incidente, masa estelar y metalicidad quedan 2\,565 planetas en 1\,981 estrellas: 848
super-Tierras y 1\,717 sub-Neptunos (66.9\,\%). Todos pertenecen también a la muestra de
cinco regímenes.

\paragraph{Análisis II.}
Es la submuestra más restrictiva. De los 1\,504 planetas con masa verdadera y radio medido
en tránsito, 1\,488 tienen además masa estelar, metalicidad y flujo incidente, y 574 de
ellos tienen entre 1 y 4~$R_\oplus$. Estos 574 planetas, en 385 estrellas, son el 9.2\,\%
del catálogo. A diferencia de la comparación anterior, esta muestra conserva los 85
planetas del intervalo de 1.6 a 2.0~$R_\oplus$, porque la respuesta es la masa y no una
clase; incluye además 120 super-Tierras y 369 sub-Neptunos. Su procedencia es distinta de
la de las otras submuestras: el 43.6\,\% de los planetas se descubrió con TESS, el 37.1\,\%
con \textit{Kepler} y el 13.6\,\% con K2. En 134 de las 385 estrellas (34.8\,\%) hay más de
un planeta de la muestra.

\paragraph{Análisis III.}
La unidad de análisis es el sistema. El catálogo tiene 3\,238 planetas de 1 a 4~$R_\oplus$
y período menor que 100~días con masa estelar y metalicidad, en 2\,340 estrellas. De
ellas, 615 tienen al menos dos de esos planetas y forman la muestra de sistemas, con
1\,513 planetas. Su composición se describe en la Sección~\ref{sec:res-analisis-III}.

\paragraph{Relación entre las submuestras.}
Las submuestras no son independientes entre sí. Además de la comparación entre super-Tierras
y sub-Neptunos, que está contenida en la muestra del análisis~I, 550 de los 574 planetas del
análisis~II y 1\,359 de los 1\,513 del análisis~III pertenecen también al análisis~I, y 312
planetas están en las cuatro. En conjunto, los tres análisis usan 4\,622 planetas distintos,
el 73.8\,\% del catálogo. Los tres análisis no equivalen, por lo tanto, a tres pruebas
independientes de la relación entre la estrella y el planeta; este punto se retoma al
compararlos (Sección~\ref{sec:res-comparacion}).

Las cuatro submuestras cubren además un intervalo estelar parecido. La mediana de la masa
estelar de sus planetas está entre 0.87 y 0.94~$M_\odot$ y la de la metalicidad, entre
$-0.01$ y 0.02~dex. En la muestra del análisis~I, el 90\,\% central de los planetas orbita
estrellas de 0.51 a 1.34~$M_\odot$ con metalicidad entre $-0.28$ y 0.31~dex. Las muestras
de los análisis~II y~III se extienden algo más hacia estrellas de baja masa (percentil 5 en
0.34~$M_\odot$). Las asociaciones que se presentan en las secciones siguientes se
refieren, entonces, a ese intervalo de estrellas.

% ----------------------------------------------------------------------------
\section{Análisis I: regímenes de radio}
\label{sec:res-analisis-I}
% Pendiente.

\subsection{Modelo multinomial de cinco regímenes}
\label{sec:res-multinomial}

\subsection{Análisis enfocado en super-Tierras y sub-Neptunos}
\label{sec:res-zoom}

\subsection{Verificaciones de robustez}
\label{sec:res-robustez-I}

% ----------------------------------------------------------------------------
\section{Análisis II: propiedades físicas de planetas pequeños}
\label{sec:res-analisis-II}
% Pendiente.

\subsection{Relación masa--radio--irradiación}
\label{sec:res-masa-radio}

\subsection{Contribución de las propiedades estelares}
\label{sec:res-contribucion-estelar}

% ----------------------------------------------------------------------------
\section{Análisis III: arquitectura de sistemas}
\label{sec:res-analisis-III}
% Pendiente.

\subsection{Espaciamiento orbital}
\label{sec:res-espaciamiento}

\subsection{Sistemas \textit{gapped}}
\label{sec:res-gapped}

% ----------------------------------------------------------------------------
\section{Comparación entre los tres análisis}
\label{sec:res-comparacion}
% Pendiente.
